\subsection{Good integrators and comparison of topologies}
\label{sec:emery-completeness}
\begin{lemma}[Locally FV processes are good integrators]
Let $A$ be a strongly adapted right-continuous real process whose paths
all have local finite variation. Under a finite measure it is an
elementary good integrator.
\end{lemma}
\begin{proof}
Let predictable elementary integrands $H_n$ tend uniformly to zero over
all times and sample points. Fix $T$ and stop $A$ there, giving paths of
global finite variation without changing the elementary gain at $T$.
Pathwise the gain has the Stieltjes representation
\[
 (H_n\cdot A)_T=\int1_{(0,T]}(t)H_n(t)\,dA^T_t.
\]
The integrands tend pointwise to zero and eventually have absolute value
at most one. The pathwise total-variation measure is finite, so dominated
convergence for the vector measure makes the gain tend to zero on each
sample path. Strong adaptation and right continuity imply progressive
measurability, hence measurability of the elementary gains. Under finite
measure their almost-sure convergence implies convergence in probability.
No bound on the number of strategy blocks or on expected variation is used.
\end{proof}
Linearity of elementary gains reduces the decomposable case to martingales,
and indistinguishability transfers the result even to nonregular raw targets.
\begin{lemma}[Martingale integral estimate]
For a right-continuous true martingale $M$ with $M_T\in L^2$, and elementary
predictable $H$ with everywhere bound $|H|\leq B$, $B\geq0$,
\[
 \|(H\cdot M)_T\|_2\leq B\sqrt{E(M_T-M_0)^2}.
\]
A right-continuous martingale square integrable at each time is consequently
a good integrator.
\end{lemma}
\begin{proof}
Stop at $T$ and take chronological factorial-grid left-endpoint integrals.
Discrete isometry bounds their squared expectation by $B^2E(M_T-M_0)^2$.
The grids cover the stopped terminal horizon, and right continuity gives
pathwise convergence to the elementary gain. Fatou for the $L^2$ norm passes
the bound to the limit. For uniformly vanishing $H_n$, write
$C=\sqrt{E(M_T-M_0)^2}$ and eventually use
$|H_n|\leq\varepsilon/(C+1)$ to get norm at most $\varepsilon$.
Measurability and $L^2$ convergence imply convergence in probability.
No block count or coefficient absolute-sum bound occurs.
\end{proof}
\begin{lemma}[Removing localization from good-integrator behavior]
If strongly adapted right-continuous $S$ under a finite measure has
$\tau_r\to\infty$ almost surely and every $S^{\tau_r}$ is a good integrator,
then so is $S$. Monotonicity of the stops is unnecessary.
\end{lemma}
\begin{proof}
For fixed $T$, bounded convergence gives $P(\tau_r\leq T)\to0$.
On $\{\tau_r>T\}$ all elementary evaluation times lie before the stop,
so the gains agree. Therefore
\[
 P(|(H_n\cdot S)_T|\geq\varepsilon)\leq P(\tau_r\leq T)
 +P(|(H_n\cdot S^{\tau_r})_T|\geq\varepsilon).
\]
First choose $r$, then send $n$ to infinity.
\end{proof}
\begin{proposition}[Decomposable processes are good integrators]
Under the usual conditions a zero-start process with admissible decomposition
$X=N+A$ is an elementary good integrator.
\end{proposition}
\begin{proof}
Apply DDY with threshold one: $N=L+Q$, $|\Delta L|\leq2$ and $Q$ locally FV.
Common localization makes $L^{\tau_r}$ a deterministically bounded true
martingale. The $L^2$ estimate and removal of localization give good-integrator
behavior for $L$. Add the locally FV $Q+A$ and transfer to the raw target by
indistinguishability. Local square integrability of $N$ itself is not assumed.
\end{proof}
\begin{lemma}[Finite stopping preserves good integrators]
If $S$ is a good integrator and $\tau$ finite-valued, so is $S^\tau$.
\end{lemma}
\begin{proof}
Stop an elementary strategy by subtracting its predictable post-stop tail.
Its integrand is $1_{\{t\leq\tau\}}H_t$, preserving uniform convergence to
zero. Commuting minima in each finite summand gives
$(H\cdot S^\tau)_T=((1_{\{t\leq\tau\}}H)\cdot S)_T$.
Transfer measurability and convergence by this identity.
\end{proof}
\begin{proposition}[Stopped decomposition of a bounded-jump good integrator]
Let $X$ be zero-start, strongly adapted, c\`adl\`ag and a good integrator,
with $|\Delta X|\leq c$ outside one common null set, $c\geq0$.
The exhaustive stops

\[
 \sigma_n=\inf\{t:|X_t|>n+1\}\wedge(n+1)
\]
 make $X^{\sigma_n}$ bounded
sources with bound $n+1+c$. Theorem~\ref{thm:bounded-source} therefore
supplies a global special decomposition for every coordinate.
\end{proposition}
\begin{proof}
Compact boundedness gives exhaustion. The pre-passage bound and stopping
jump give the closed bound. Adaptation and c\`adl\`ag regularity survive
stopping, as does good-integrator behavior by the preceding lemma.
These supply all inputs to Theorem~\ref{thm:bounded-source}.
\end{proof}
\begin{proposition}[Removing large jumps from a general good integrator]
For zero-start strongly adapted c\`adl\`ag good integrator $X$ and $c>0$, set
\[
 J_t=\sum_{0<s\leq t}\Delta X_s1_{\{|\Delta X_s|>c\}},\qquad Y=X-J.
\]
Then $J$ is zero-start strongly adapted c\`adl\`ag locally FV, and $Y$ is
such a good integrator with $|\Delta Y_t|\leq c$ at every path and time.
Thus the preceding stopped-decomposition proposition applies to $Y$.
\end{proposition}
\begin{proof}
The finite-horizon sums $J^T$ agree before smaller horizons.
Define $J_t=J^t_t$; fixed-time adaptation follows from that coordinate,
right continuity from agreement on a larger horizon, and left limits and
local BV from compact-interval agreement. At $t>0$ the residual agrees,
including its left limit, with the finite-horizon residual at horizon $t$,
so inherits the jump bound; at zero use the convention of zero jump.
Subtract the locally FV good integrator $J$ from $X$, using linearity and
stability of convergence in probability. No compensation of $J$ is needed,
since the final FV component need only be adapted.
\end{proof}
\begin{theorem}[Decomposition of regular zero-start good integrators]\label{thm:regular-gi-decomposition}
Under the usual conditions, a zero-start strongly adapted c\`adl\`ag $X$
is an elementary good integrator if and only if it has a decomposition
$X=N+B$ with zero-start c\`adl\`ag local martingale $N$ and zero-start
strongly adapted c\`adl\`ag locally FV $B$.
\end{theorem}
\begin{proof}
One direction was proved above. For the other remove jumps larger than one,
writing $X=Y+J$. Decompose the bounded stopped sources
$Y^{\sigma_n}=M^n+A^n$ and normalize
\[
 \widehat M^n=(M^n-M^n_0)^{\sigma_n},\qquad
 \widehat A^n=(A^n-A^n_0)^{\sigma_n}.
\]
Local martingales can be centered without global integrability of the initial
value. Constant extension of a predictable initial value and finite stopping
preserve predictability. Each stopped FV path has global bounded variation.
Zero start of the source and idempotence of stopping preserve the decomposition.
Stop two coordinates at $\rho=\sigma_n\wedge\sigma_m$.
Their FV difference is indistinguishable from a local-martingale difference,
and is zero-start predictable right-continuous globally BV. Rigidity makes
it zero, giving compatibility of both components.
Glue the same martingale and predictable FV families. Exhaustion gives
$Y=M+A$ on a common full-measure set. The martingale component is strongly adapted and c\`adl\`ag; the other
component is predictable, right-continuous and locally of finite variation.
Local finite variation gives its left limits. Recenter after gluing and set
\[
 N=M-M_0,\qquad B=A-A_0+J.
\]
Since $M_0+A_0=0$ almost surely, $X=N+B$.
The final $B$ is adapted locally FV; predictability is not asserted.
\end{proof}

\begin{proposition}[Passing uniform Cauchy estimates to a ucp limit]
Let $X_n$ be strongly progressively measurable and right-continuous, and
let $X_n\to Y$ ucp. For a unit-bounded predictable elementary test $J$, put
\[
 \bar e_T^J(U,V)=E[1\wedge\sup_{t\leq T}|(J\cdot U)_t-(J\cdot V)_t|].
\]
Suppose that, for every $T,\varepsilon>0$, there is $N$ such that
$\bar e_T^J(X_m,X_n)\leq\varepsilon$ for all $m,n\geq N$ and all $J$.
Then $\sup_J\bar e_T^J(X_n,Y)\to0$. Conversely, this uniform convergence implies
the Cauchy condition.
\end{proposition}
\begin{proof}
For fixed $J$, the finite-sum estimate bounds the capped maximal test error
by the capped maximal process error multiplied by
\[
 \max\{8\operatorname{length}(J),1\}.
\] This constant is used only for a fixed test. The ucp convergence
and boundedness of capped errors in $[0,1]$ give $\bar e_T^J(X_m,Y)\to0$.

For $T,\varepsilon$, choose the Cauchy cutoff $N$ first, then fix $n\geq N$
and $J$. For all $m\geq N$, the triangle inequality gives
\[
 \bar e_T^J(X_n,Y)\leq \bar e_T^J(X_m,X_n)+\bar e_T^J(X_m,Y)
 \leq\varepsilon+\bar e_T^J(X_m,Y).
\]
Send $m\to\infty$ to get $\bar e_T^J(X_n,Y)\leq\varepsilon$.
Only now take the supremum over $J$, since $N$ was chosen independently
of the test. The converse uses the triangle inequality through $Y$ and
an error bound $\varepsilon/2$ for each of its two terms.
\end{proof}
No closedness of a fixed source's integral range is used. We construct the
regular limit and its good-integrator behavior next.
\begin{proposition}[A fast subsequence on increasing horizons]
Suppose strongly progressively measurable right-continuous processes $X_n$
satisfy the uniform elementary-test Cauchy condition and
$X_n(0)=X_0(0)$ almost surely for every $n$. There is a strictly increasing
$f$ such that, almost surely, for all sufficiently large $k$,
\[
 \sup_{t\leq k+1}|X_{f(k+1)}(t)-X_{f(k)}(t)|\leq2^{-k-1}.
\]
Thus on each finite interval the sufficiently late adjacent suprema are
summable pathwise.
\end{proposition}
\begin{proof}
The unit test $1_{(0,T]}$ integrates to $X(t)-X(0)$ for $t\leq T$.
The common initial value therefore turns the uniform test estimate into
\[
 E[1\wedge\sup_{t\leq T}|X_m(t)-X_n(t)|]\leq\varepsilon
\]
for all sufficiently large $m,n$. Put $q_k=2^{-k-1}$ and let $N_k$ be the
cutoff at $T=k+1$ and error $q_k^2$. Set $f(0)=N_0$ and
$f(k+1)=\max(f(k)+1,N_{k+1})$.
For the consecutive capped maximum $E_k$, we have $EE_k\leq q_k^2$.
Markov gives $P(E_k\geq q_k)\leq q_k$.
Since $\sum_kq_k<\infty$, Borel--Cantelli gives $E_k<q_k$ eventually,
outside one null set. As $q_k<1$, the cap can be removed.
Right continuity makes the stopped factorial grid detect the entire interval
supremum, so the estimate holds simultaneously for all $t\leq k+1$.
\end{proof}
\begin{proposition}[Regular limits of test-uniform Cauchy sequences]
Assume the usual conditions. If strongly progressively measurable
c\`adl\`ag processes $X_n$ are Cauchy uniformly over elementary tests and $X_n(0)=X_0(0)$ almost surely for every $n$, there is a
strongly progressively measurable c\`adl\`ag process $Y$ such that
\[
 Y(0)=X_0(0)\quad\text{a.s.},\qquad
 \sup_J\bar e_T^J(X_n,Y)\longrightarrow0\quad(T<\infty).
\]
The entire sequence also converges to $Y$ ucp.
\end{proposition}
\begin{proof}
Take the preceding fast subsequence. On each finite interval, outside one
common null set, sufficiently late adjacent suprema are bounded by a
summable numerical sequence. The triangle inequality and series-tail
bound give locally uniform Cauchy paths. Completeness of the reals gives
a limit $Z(t)$ at every time, with locally uniform convergence.
Use a measurable limit selection defined even on nonconvergent sample
points. Then $Z(t)$ is measurable with respect to the information at $t$,
so $Z$ is strongly adapted.

Stopping paths at a finite horizon reduces this to all-time uniform
convergence. Uniform limits preserve right continuity and left limits,
so $Z$ is almost surely c\`adl\`ag. The usual conditions put the exceptional
set in the time-zero sigma algebra. Replacing its paths by zero gives
strongly adapted, everywhere c\`adl\`ag $Y$. Adapted right continuity
implies strong progressive measurability. Locally uniform convergence
survives outside the common null set, and convergence at time zero gives
the asserted initial value.

The subsequence's local uniform convergence gives convergence in probability
of every capped maximal error. In the Cauchy-to-limit proof let the
auxiliary index $m$ run through $f(m)$. Since $f(m)\to\infty$, the same
Cauchy cutoff applies. Take the fixed-test limit first to obtain uniform
error bounds for the original entire sequence. Thus that sequence
converges against elementary tests. The common initial value and unit
test give full-sequence ucp convergence as well.
\end{proof}
This construction does not assume the approximants are good integrators or
use integral-range closedness.
\begin{proposition}[Good-integrator stability under test-uniform limits]\label{prop:gi-limit}
Suppose strongly progressively measurable right-continuous $X_n,Y$ satisfy
$\sup_J\bar e_T^J(X_n,Y)\to0$ for each finite $T$. If every $X_n$ is a good
integrator, so is $Y$.
\end{proposition}
\begin{proof}
Fix uniformly vanishing predictable elementary integrands $H_k$ and a
horizon $T$. For arbitrary $\varepsilon,\delta>0$, put
$a=\min(\varepsilon/2,1)>0$.
Choose one $m$ such that $\bar e_T^J(X_m,Y)\leq a\delta/4$ for every unit test.
With this $m$ fixed, good-integrator behavior of $X_m$ gives, eventually,
\[
 P(|(H_k\cdot X_m)_T|\geq\varepsilon/2)<\delta/2.
\]
On the same tail $|H_k|\leq1$. The capped maximal integral difference
$E_k$ therefore has $EE_k\leq a\delta/4$, so Markov yields
$P(E_k\geq a)\leq\delta/4$.
Right continuity bounds the capped difference at $T$ by $E_k$.
The triangle inequality gives the event inclusion
\[
 \{|(H_k\cdot Y)_T|\geq\varepsilon\}
 \subseteq\{|(H_k\cdot X_m)_T|\geq\varepsilon/2\}\cup\{E_k\geq a\}.
\]
Its probability is less than $\delta$. Since $\delta$ was arbitrary,
this proves convergence in probability. Measurability of the gain follows
from strong progressive measurability of $Y$ and its finite-sum formula.
\end{proof}
\begin{corollary}[The Cauchy limit remains decomposable]
Under the usual and filtration sigma-finiteness conditions, a test-uniform
Cauchy sequence of regular zero-start good integrators has a regular
zero-start decomposable limit, with full-sequence elementary-test convergence.
\end{corollary}
\begin{proof}
Construct $Y$ as above and set $Z_t=Y_t-Y_0$.
Centering preserves all elementary increments, regularity, and adaptation.
Stability makes $Z$ a good integrator; apply its zero-start decomposition
theorem. The FV component need only be adapted.
\end{proof}
\begin{proposition}[Elementary gauges on the same quotient]
On the indistinguishability quotient $\mathcal S$ of decomposable zero-start
processes,
\[
 e_T([X],[Y])=\sup_J\bar e_T^J(X,Y)
\]
 is well-defined.
Every sequence Cauchy for all finite-horizon gauges has a unique gauge limit
in $\mathcal S$.
\end{proposition}
\begin{proof}
Indistinguishable sources have identical finite-sum test integrals at every
time off one null set. Envelopes, integrals, and test suprema are unchanged.
Raw representatives need not be regular: choose each decomposition's component
sum, which is zero-start strongly progressive and c\`adl\`ag.
This preserves the Cauchy condition. Decomposability gives good-integrator
behavior, so the preceding corollary constructs a limit in the same domain.
Return to the original representatives and take the quotient.
For uniqueness use regular zero-start versions, fixed-test uniqueness and
the unit test to identify increments, then initial values to identify paths.
\end{proof}
This construction uses no auxiliary-metric completeness or continuity,
only identification of its zero-distance relation with indistinguishability.
\begin{proposition}[A time-weighted elementary metric]
On the same quotient,

\[
 d_{\mathcal E}(x,y)=\sum_{n\geq0}2^{-n-1}e_{n+1}(x,y)
\]
 is a metric bounded
by one.
\end{proposition}
\begin{proof}
Regular representatives make test errors integrable and in $[0,1]$.
Symmetry and triangle pass through integration and the test supremum;
horizon monotonicity follows from the all-time supremum on each interval.
The weights sum to one, so the series is finite and has the metric algebraic
properties. If its value is zero, each positive-weight gauge is zero.
Horizon monotonicity gives zero at every finite horizon; uniqueness of the
gauge limit of a constant sequence gives equality of the two quotient points.
\end{proof}
\begin{proposition}[Independent completeness of the elementary metric]
$(\mathcal S,d_{\mathcal E})$ is complete.
\end{proposition}
\begin{proof}
For $T\leq k+1$,
\[
 2^{-k-1}e_T(x_m,x_n)\leq d_{\mathcal E}(x_m,x_n).
\]
A metric-Cauchy sequence is therefore Cauchy for every gauge.
Take its already constructed gauge limit and use dominated convergence
with majorant $2^{-k-1}$ to obtain metric convergence.
Neither auxiliary completeness nor closedness of an integral range is used.
\end{proof}
\begin{proposition}[The elementary metric and addition]
The natural real-linear structure on $\mathcal S$ has translation-invariant
$d_{\mathcal E}$, and subtraction is $2$-Lipschitz for the maximum product metric.
\end{proposition}
\begin{proof}
Linearity in the source makes each test difference invariant under adding
the same process. Integrate, take suprema, and sum the time weights.
Symmetry also gives invariance under simultaneous negation. Thus
\[
 d_{\mathcal E}(x-y,x'-y')\leq d_{\mathcal E}(x,x')+d_{\mathcal E}(y,y')
 \leq2\max(d_{\mathcal E}(x,x'),d_{\mathcal E}(y,y')).
\]
This gives the compatible additive uniform structure and continuous addition
and negation.
\end{proof}
\begin{proposition}[Fixed scalar multiplication]

\[
 d_{\mathcal E}(ax,ay)\leq\max(|a|,1)d_{\mathcal E}(x,y)
\]
 for every real $a$.
\end{proposition}
\begin{proof}
For positive $a$ use finite-sum linearity and
$1\wedge(au)\leq\max(a,1)(1\wedge u)$, then integrate, take suprema, and
sum. Zero is immediate; negative $a$ reduces to positive by negation isometry.
\end{proof}
The factor does not tend to zero with $a$; the next estimate handles that case.
\begin{lemma}[Uniform boundedness of unit-integral maxima]
For a strongly adapted right-continuous good integrator $S$, right-continuous
filtration, and finite $T$, the maxima $\sup_{t\leq T}|(H\cdot S)_t|$ over
all unit-bounded elementary predictable $H$ are bounded in probability.
\end{lemma}
\begin{proof}
The good-integrator property gives boundedness of unit integrals at the
fixed time $T$. Choose a tail bound $r\geq0$ and put $R=r+1$.
Let $\tau$ be the first strict absolute passage of $X=H\cdot S$ above $R$.
It is a stopping time, and if finite then $|X_\tau|\geq R$ by right continuity.
Stopping $H$ at $\tau\wedge T$ gives another unit strategy $K$.
If the original maximum exceeds $R$, then $\tau\leq T$ and
$|(K\cdot S)_T|\geq R>r$. Its fixed-time tail controls the original maximal tail.
\end{proof}
\begin{proposition}[Joint scalar continuity for the elementary metric]
$(\mathcal S,d_{\mathcal E})$ is a complete real topological vector space.
\end{proposition}
\begin{proof}
Fix $x$ and a regular representative $S$, which is a good integrator.
At horizon $T$ choose $R>0$ making all unit-integral maximal tails at most
$\varepsilon/2$. If $|a|\leq\varepsilon/(4R)$, the capped scaled maximum $f$
is at most $\varepsilon/4$ on the event where the unscaled maximum is at
most $R$. Hence $P(f\geq\varepsilon/2)\leq\varepsilon/2$, and
$f\leq\varepsilon/2+1_{\{f\geq\varepsilon/2\}}$ gives $Ef\leq\varepsilon$.
Here the random variable being bounded is
\[
 f=1\wedge\sup_{t\leq T}|a(H\cdot S)_t|.
\]
The scalar cutoff is test-independent, so $e_T(ax,0)\to0$, and therefore
the weighted metric tends to zero.
For $a_n\to a,x_n\to x$, use
\[
 d_{\mathcal E}(a_nx_n,ax)\leq
 \max(|a_n|,1)d_{\mathcal E}(x_n,x)+d_{\mathcal E}((a_n-a)x,0)\to0.
\]
Sequential continuity is continuity in metric spaces. Addition and
completeness were already established.
\end{proof}
\begin{proposition}[The auxiliary complete topological vector space]
The same quotient with auxiliary metric $R$ and its natural real-linear
operations is a complete metrizable real topological vector space.
\end{proposition}
\begin{proof}
Admissible decompositions are stable under addition, negation, and scaling.
The triangle inequality gives continuity of addition and
\[
 R((X-Y)-(X'-Y'))\leq R(X-X')+R(Y-Y')
 \leq2\max(R(X-X'),R(Y-Y')),
\]
so subtraction is uniformly continuous. Use the earlier scalar-continuity
lemma. The operations descend to the separated quotient and its quotient
map is continuous linear. Its equality is indistinguishability, and the
already proved completeness is for this same metric.
\end{proof}
\begin{proposition}[Closed graph between the two topologies]
The identity $(\mathcal S,d_{\mathcal E})\to(\mathcal S,R)$ has closed graph.
\end{proposition}
\begin{proof}
Suppose $x_n\to y$ in the elementary metric and $x_n\to z$ in the auxiliary
metric. Choose regular zero-start representatives from decompositions.
Positive time weights and horizon monotonicity turn elementary-metric
convergence into every finite-horizon gauge convergence. The unit test and
common initial values give ucp convergence to $y$. The auxiliary maximal-tail
estimate gives ucp convergence to $z$. Uniqueness of right-continuous ucp
limits makes $y,z$ indistinguishable. The product is metrizable, so sequential
closedness proves closedness of the graph.
\end{proof}
\begin{lemma}[Open mapping and closed graph for complete invariant metrics]
Let $E,G$ be real topological vector spaces with translation-invariant metrics.
If $E$ is complete and $G$ is Baire, every continuous surjective linear map
$f:E\to G$ is open. If both are complete, a linear map with closed graph
is continuous.
\end{lemma}
\begin{proof}
A zero neighborhood $V$ absorbs $E$, giving
$G=\bigcup_{n\geq0}(n+1)\overline{f(V)}$.
Baire and the homeomorphism given by nonzero scaling make
$\overline{f(V)}$ have nonempty interior.
For any zero neighborhood $U$ choose symmetric $V$ with $V+V\subset U$;
subtract an interior point to see that $\overline{f(U)}$ is a zero neighborhood.
For $\varepsilon>0$ choose
\[
 0<\delta_n\leq2^{-n},\qquad
 B_G(0,\delta_n)\subset
 \overline{f(B_E(0,(\varepsilon/4)2^{-n}))}.
\]
For $d_G(y,0)<\delta_0$, construct successive corrections from $x_0=0$
with $d_G(f(x_n),y)<\delta_n$ and
$d_E(x_n,x_{n+1})\leq(\varepsilon/4)2^{-n}$.
Completeness gives $x_n\to x$ with $d_E(x,0)\leq\varepsilon/2$;
continuity and vanishing residual give $f(x)=y$. Thus the image of the
$\varepsilon$-ball contains a zero neighborhood. Translate to obtain openness.
A closed graph is a complete closed linear subspace of the product.
Apply openness to its first projection; its continuous inverse followed by
the second projection is the original map.
\end{proof}
\begin{proposition}[From elementary convergence to auxiliary control]\label{prop:topology-localization}
For decomposable processes, elementary-test convergence $X_n\to X$ implies
$R(X_n-X)\to0$. For elementary representatives $G_n$ of a realized gain $G$
over a strongly progressive c\`adl\`ag good-integrator source, one subsequence
$n_k$ and one exhaustive $\tau_r\leq r+1$ give prefix witnesses for
$G_{n_k}-G$ with costs $c_{r,k}$ satisfying
\[
 \sum_kc_{r,k}\leq6(r+4)<\infty,\qquad c_{r,k}\longrightarrow0.
\]
A simultaneous selection for consecutive differences has bound $6(r+5)$.
Assume the usual conditions and sigma-finite filtration.
\end{proposition}
\begin{proof}
The closed-graph theorem makes the identity from elementary to auxiliary
topology continuous. Every elementary gain has a regular zero-start
good-integrator decomposition. Independent elementary completeness constructs
a decomposable limit, identified with the designated gain by ucp uniqueness.
Apply the auxiliary subsequence and common-localization theorem.
For consecutive differences first select
$\sum_kR(G_{n_k}-G)\leq1$. Triangle gives a difference-cost sum at most $2$;
apply summable-cost localization without another extraction.
Realization of the limit as an integral has not been assumed.
\end{proof}
This comparison is independent of closedness of a fixed source's integral range.
\begin{proposition}[One measure for prefix costs and boundary jumps]\label{prop:joint-costs}
For an everywhere bounded source, one can choose simultaneously an equivalent
probability $Q$, a subsequence $n_k$, and exhaustive $\tau_r\leq r+1$ so
that for $Z_k=G_{n_{k+1}}-G_{n_k}$,
\[
 \sum_kc_{r,k}\leq6(r+5),\qquad
 \sum_kE_Q|\Delta Z_k(\tau_r)|<\infty.
\]
There is also a measurable $\eta\geq0$ in $L^2(Q)$, defining $Q$ by an
exponential tilt with bounded density, which bounds $|G_{n_k}-G|$ on
$[0,k+1]$ simultaneously almost surely. Thus $\xi\in L^2(\mu)$ implies
$\xi+\eta\in L^2(Q)$.
\end{proposition}
\begin{proof}
Fix $Q$ first from growing-horizon error envelopes. Both auxiliary costs
and these envelopes' $L^1(Q)$ norms vanish; choose a subsequence whose sum
is bounded termwise by a geometric sequence. Increasing indices allow
shrinking the measured horizon to $k+1$, retaining the envelope and sum bound
one. Consecutive auxiliary costs sum to at most two, giving common stops.
For fixed $r$ and $k+1\geq r+1$, if $a_k$ is the error envelope then

\[
 |\Delta Z_k(\tau_r)|\leq2(a_{k+1}+a_k),
\] with summable expectations.
The finitely many earlier terms are integrable because the bounded source
and each row's finite coefficient bound bound its entire gain.
No further extraction changes the common stops.
\end{proof}

\paragraph{Comparing prefix witnesses with actual component costs.}
For the same witnesses define
$\widehat N_k=(N_k-N_k(0))^{\tau_r}$.
Intrinsic quadratic variation and reverse Davis give
\[
 E_Q\sqrt{[\widehat N_k]_{r+1}}
 \leq7E_Q\sup_{t\leq r+1}|\widehat N_k(t)|\leq14c_{r,k}.
\]
Centering costs at most twice the original maximum; neither bounded jumps
nor a separately supplied quadratic schedule is needed.
For a prefix witness $N,A$, set
\[
 M=((N-N_0)^\tau)^T,\qquad
 B=A^{\tau-}-A_0-\Delta N_\tau1_{[\tau,\infty)}.
\]
Use the zero-jump convention at $\tau=0$.
The second component is adapted, zero-start, and globally FV. Their sum
agrees with the zero-start target before $\tau$; when $\tau=0$ that condition
is empty. Prefix variation and the boundary jump give $E\Var(B)\leq2c$,
while the martingale root expectation is at most $14c$.
Thus the actual J1 decomposition bounds the strict-prefix infimum,
here denoted $J_\tau(X)$, as follows:
\[
 J_\tau(X)\leq16c.
\]
Indistinguishability of the exact witness transfers initial values and prefix
agreement on the same null set. Consequently the same rows have summable
$J_{\tau_r}(Z_k)$.
If an integral representation is supplied, Lemma~\ref{lem:canonical-cost}
and equation~\eqref{eq:actual-closed-cost} give
\[
 E\sup_{t\leq T}|M^{\mathrm{actual}}_t-M^{\mathrm{actual}}_0|
 +E\Var_{[0,T]}A^{\mathrm{actual}}
 \leq30\bigl(32c+E|\Delta X_\tau|\bigr).
\]
This estimate does not itself construct the integral representation and does
not transfer special certificates unconditionally between measures.

\subsubsection{Constructing actual rows from the same selection.}\label{sec:actual-series}
Choose the preceding $Q,f,\tau_r$ from the bounded source and elementary
representatives. An almost-sure bound uniform over time suffices: the finitely
many initial jump terms can be bounded off that common null set, and the
bound transfers by $Q\ll\mu$.
After fixing $Q$, construct the original price's decomposition and completed
schedule. Each $Z_k$ is the gain of the difference of two elementary strategies,
whose coefficient absolute-sum bound is the sum of their row bounds.
Construct its stopped graph directly at $\tau_r$, and stop again at
$T_r=r+1$. The gain is still $Z_k^{\tau_r}$ and the coefficient is
\[
 1_{(0,\tau_r]}(H^{f(k+1)}-H^{f(k)}).
\]
Both components are constant after $T_r$ and the FV component is globally BV.
These are integral certificates for the original $S$.
Indistinguishability of the stopped and represented gains allows the
comparison estimate to apply to these independently constructed certificates.
For their components $M^{r,k},A^{r,k}$ it yields
\[
 \sum_k\left(E_Q\sup_{t\leq T_r}|M^{r,k}_t-M^{r,k}_0|
 +E_Q\Var_{[0,T_r]}A^{r,k}\right)<\infty.
\]
The tilt, density bound, and $L^2(Q)$ error envelope remain part of the same
selection.

\paragraph{All-time component series.}
Fix $r$ and center $m_k=M^{r,k}-M^{r,k}_0$, $a_k=A^{r,k}-A^{r,k}_0$.
The actual rows are restopped elementary realizations retaining martingale
left limits. Each $m_k$ is c\`adl\`ag and constant after $T_r$, so its
measurable horizon envelope $g_k$ dominates it for all time, with
$\sum_kE_Qg_k<\infty$.
This upgrades each local martingale to a true martingale; the dominated
series theorem makes
\[
 m(t)=\sum_km_k(t)
\]
 a true martingale.
Partial sums converge uniformly over all time almost surely.
Centering does not change FV, and constancy after $T_r$ gives
\[
 \Var_{[0,\infty)}a_k\leq\Var_{[0,T_r]}A^{r,k}.
\]
Tonelli makes the variation series finite almost surely, giving a globally
BV sum $a(t)=\sum_ka_k(t)$ with all-time uniform and variation-distance
convergence. It is predictable as a pointwise series of predictable components.
All conclusions concern the actually generated rows of the same selection.

\paragraph{Actual approximants including the first term.}
Write $G_k=H^{f(k)}\cdot S$. At stage $r$ let $P_{r,0}$ be the directly
constructed certificate of $G_r^{\tau_r}$ and recursively, in the raw layer,

\[
 P_{r,n+1}=P_{r,n}+R_{r,r+n}
\]
 using the same actual increment rows.
Telescoping gives coefficient $1_{(0,\tau_r]}H^{f(r+n)}$ and gain
$G_{r+n}^{\tau_r}$. Transfer realization from its direct elementary certificate
using coefficient equality and gain indistinguishability; a general actual
addition calculus is unnecessary. Initial and increment components are zero-start.
If $c_k$ bounds the $k$th elementary gain, the simultaneous error envelope gives
\[
 |G_{r+n}^{\tau_r}(t)|\leq |G_r^{\tau_r}(t)|
 +|(G_{r+n}-X)^{\tau_r}(t)|+|(G_r-X)^{\tau_r}(t)|
 \leq c_r+2\eta=:\zeta_r.
\]
Here $\tau_r\leq r+1\leq r+n+1$. Thus $\zeta_r\in L^2(Q)$ is a common
all-time gain envelope for these actual approximants. Neither a new measure
nor an additional $L^2$ assumption on the designated $X$ is needed.

\paragraph{Regular component limits and the Mémín input.}
Right continuity and left limits pass through all-time uniform convergence.
Correct bad paths on time-zero measurable null sets, intersecting the two
component agreement events. Preserve the martingale property, predictability,
uniform convergence, and variation-distance convergence.
Add the initial components to obtain $M_r^\infty,A_r^\infty$; the former is
an everywhere c\`adl\`ag local martingale and the latter an everywhere
right-continuous predictable globally BV process. Global BV is asserted at
this fixed stopped stage only.
Adjacent raw finite sums differ by exactly the stored increment components.
Because the FV paths are constant after $T_r$, their canonical Stieltjes
measure's total variation equals the path variation on $[0,T_r]$.
This uses the stopped-path theorem, not an auxiliary measure field in a record.
Every horizon martingale-difference envelope is likewise bounded by the
$T_r$ envelope. The expected-cost sum and Tonelli therefore make both
adjacent-difference norm series finite almost surely.
Retain the stronger selected error bound
\[
 \sum_kE_Q\sup_{t\leq k+1}|G_k(t)-X(t)|\leq1.
\]
Tonelli makes the envelopes summable pathwise and tending to zero.
On one full-measure set $G_k(t)\to X(t)$ at all times, hence also at
$t\wedge\tau_r$. Gain identification and component convergence give
\[
 M_r^\infty+A_r^\infty\sim X^{\tau_r}.
\]
Together with $\zeta_r$ these supply every Mémín component-data field and
the stored stopped-gain identification from the source and representatives.
No new subsequence or measure is selected.

\paragraph{Component identification by stopping and finite addition.}
For a right-continuous local martingale $M$ with localizers $\rho_n$, put
$B_n=\{\rho_n>0\}$. The process $1_{B_n}M$ stopped at $\rho_n$ is a
martingale. Stop further at $\sigma$ and commute the indicator and the two
stops. The same localizers establish the local property of $M^\sigma$;
this argument does not require $M_0=0$.
Finite stopping preserves predictability of a predictable FV component and
turns local BV into global BV. Thus stop raw gains and both components
directly, keeping initial values. On a graph coordinate the raw stop and
the completed strategy have the same gain; uniqueness of centered canonical
components identifies both parts. Repeat on minimum stops of two graph
schedules, identifying energy kernels, variation bridges, and completed
integrals. This supplies finite addition on a common schedule.
Together with the finite-horizon Stieltjes identifications this constructs
the needed FV calculus without assuming a general predictable-restriction
calculus. The intrinsic special terminal model using this addition does not
require a stopping calculus as input; identifying it with the general market
is a separate step.

\subsubsection{Direct joint passages and realization of the stopped limit.}\label{sec:joint-realization}
At fixed outer stage $r$ combine the source completed schedule with all
martingale passages of $P_{r,k}$ to obtain exhaustive finite-horizon stops
$\rho_{r,j}$. Construct source prefixes by the direct stopping rule.
The pre-passage bound and the Corollary~2.4 jump envelope include overshoot
and give
\[
 \|M_{r,k}^{\rho_{r,j}}(T_j)\|_{L^2(Q)}
 \leq c_j+6\|\zeta_r\|_{L^2(Q)}.
\]
The stopped local martingale has an integrable all-time envelope, so is a
true martingale by dominated convergence.
The fixed graph's completed coordinate and the directly stopped raw
components have the same gain. Centered uniqueness identifies the integral
representation of each approximating component. Apply
Lemma~\ref{lem:joint-coefficient} with the summable adjacent FV variations
and the common-coordinate $L^2$ bounds on martingale terminals proved above.
That lemma first obtains convergence of the full coefficient sequence under
variation control, then makes forward averages of the same coefficients
converge in the energy coordinates. One predictable coefficient consequently
represents both limiting integrals. Agreement with the variation-control
limit also determines the FV integral on energy-null regions.
The exhaustive coordinate graph realizes the component-sum limit as an
actual strategy. Applying this to the same $Q,\tau_r$ component data gives
an original-price actual strategy $V_r$ with
\[
 (V_r\cdot S)\sim_Q X^{\tau_r},\qquad
 \tau_r\leq r+1,\qquad\tau_r\to\infty\quad Q\text{-a.s.}
\]
The left side denotes its stored gain. No new stopping calculus, measure,
or approximating sequence is an input. These are special certificates at
each stopped stage under $Q$. Return to the original measure's general graph,
pasting of stopping intervals, and terminal identification are the separate
steps already given in Section~\ref{sec:integral-calculus}.
